% b3_trinom.tex -- DRAFT, local file only. Not submitted, posted or shared.
\documentclass[11pt]{amsart}
\usepackage{amsmath,amssymb,amsthm}
\usepackage{url}
\usepackage[colorlinks=true,linkcolor=blue,citecolor=blue,urlcolor=blue]{hyperref}
% Uzun fayl nomlari (\texttt) qator chetidan chiqmasin: bo'shliq cho'zilsin, kerak bo'lsa sloppy
\setlength{\emergencystretch}{3em}
\sloppy

\newtheorem{theorem}{Theorem}[section]
\newtheorem{proposition}[theorem]{Proposition}
\newtheorem{lemma}[theorem]{Lemma}
\newtheorem{corollary}[theorem]{Corollary}
\newtheorem{conjecture}[theorem]{Conjecture}
\theoremstyle{definition}
\newtheorem{definition}[theorem]{Definition}
\theoremstyle{remark}
\newtheorem{remark}[theorem]{Remark}

\newcommand{\Tr}{\mathrm{T}}
\newcommand{\Cz}{C_0}
\newcommand{\Ck}{C_k}
\newcommand{\F}{\mathcal{F}}
\newcommand{\TODO}[1]{\textnormal{\textbf{[TODO: #1]}}}
\newcommand{\ZenodoDOI}{\href{https://doi.org/10.5281/zenodo.23311094}{10.5281/zenodo.23311094}}

\title[The trinomials $1+ax^3+x^k$]{The trinomials $1+ax^3+x^k$ are not unfair factors\\ of $0$--$1$ polynomials}
% Author data as used in the author's RNT submission (DAVOM-RNT-PROMPT-2026-10-10).
\author{Otakhon U. Kenjaev}
\address{Independent researcher, Khorezm, Uzbekistan}
\email{kenjayevotaxon@gmail.com}
\thanks{ORCID: 0009-0009-3566-9285.}
\date{12 October 2026}

\begin{document}

\begin{abstract}
The unfair $0$--$1$ polynomial conjecture asserts that if a polynomial with
coefficients in $\{0,1\}$ factors as a product of two monic polynomials with
nonnegative real coefficients, then both factors are $0$--$1$ polynomials.
We prove that for every $k\ge 4$ with $k\not\equiv 0 \pmod 3$ and every
$0<a<1$, the trinomial $1+ax^3+x^k$ does not divide any $0$--$1$ polynomial with a
nonzero cofactor having nonnegative coefficients; by reciprocity the same holds
for $1+ax^{k-3}+x^k$. The proof analyses the forced recurrence
$q_n+aq_{n-3}+q_{n-k}=r_n\in\{0,1\}$. An exact class formula for the
coefficients along a distinguished ``main branch'' and elementary deviation
lemmas reduce the problem to the main branch; for $a\in(1/2,1)$ this is settled
by hand, and for $a\in(0,1/2]$ and $k\ge 67$ by a single rigorous
interval-arithmetic certificate that is uniform in $k$. The remaining range
$k\le 149$ is covered by exhaustive computer certificates with exact algebraic
treatment of isolated parameters.
\end{abstract}

\maketitle

% ---------------------------------------------------------------------------
\section{Introduction}

A polynomial is a \emph{$0$--$1$ polynomial} if all its coefficients lie in
$\{0,1\}$. The following question was raised on MathOverflow~\cite{MO339137}.

\begin{conjecture}[Unfair $0$--$1$ polynomial conjecture]
Let $P,Q\in\mathbb{R}[x]$ be monic with nonnegative coefficients. If $PQ$ is a
$0$--$1$ polynomial, then $P$ and $Q$ are $0$--$1$ polynomials.
\end{conjecture}

Following the literature, we call a polynomial $P$ with nonnegative coefficients
an \emph{unfair factor} if it is not a $0$--$1$ polynomial and $PQ$ is a $0$--$1$
polynomial for some nonzero $Q$ with nonnegative coefficients. Ghidelli~\cite{Ghidelli2022}
showed that $1+ax^2+x^5$ ($0<a<1$) is not an unfair factor. Hare~\cite{Hare2023}
gave an algorithm that rules out all candidate factors of degree at most $15$
with the exception of $975$ candidates (out of $7141686$ cases); the list of exceptions
does not appear to be publicly available. Hare also observed that a
computation up to index $100$ suffices to show that $1+ax^3+x^5$ is not an unfair
factor; the case $(b,k)=(3,5)$ of our theorem is therefore due to Hare. Hare's
Theorem~1.1~\cite{Hare2023} is the reciprocity principle used below.
Dvorsky~\cite{Dvorsky2026} treats the family $1+ax^2+x^k$ for odd $k\ge 341$,
the bounded odd range $7\le k\le 339$ being taken from a companion
manuscript~\cite{DvorskyCompanion}; that manuscript is not public and we have not seen
it (its metadata are as cited in~\cite{Dvorsky2026}). On the computational side, Zhang reports a verification
of the conjecture for $\deg(PQ)\le 66$~\cite{ZhangMO} (the README of the accompanying
repository~\cite{ZhangGitHub} states $\deg(PQ)\le 45$). For related work on an
$\varepsilon$-version of the problem see~\cite{NguyenPhamPham2026}.

Our main result concerns the next family of trinomials.

\begin{theorem}\label{thm:main}
Let $k\ge 4$ with $k\not\equiv 0\pmod 3$ and let $0<a<1$. Then
$P(x)=1+ax^3+x^k$ does not divide any $0$--$1$ polynomial $R$ with $R=PQ$, where
$Q\neq 0$ has nonnegative real coefficients.
\end{theorem}

\begin{corollary}\label{cor:recip}
Under the same hypotheses, $1+ax^{k-3}+x^k$ is not an unfair factor.
\end{corollary}

\begin{proof}
If $PQ=R$ then $P^*Q^*=R^*$, where $f^*(x)=x^{\deg f}f(1/x)$; $Q^*$ has
nonnegative coefficients and $R^*$ is a $0$--$1$ polynomial. Since
$P^*=1+ax^{k-3}+x^k$, the claim follows from Theorem~\ref{thm:main}. This is
Hare's Theorem~1.1~\cite{Hare2023}.
\end{proof}

\begin{remark}
If $k\equiv 0\pmod 3$ the statement is trivial; see Remark~\ref{rem:trivial}.
The families $(2,k)\cong(k-2,k)$ of~\cite{Dvorsky2026} and $(3,k)\cong(k-3,k)$
meet only for $k=5$, the case of~\cite{Ghidelli2022}. Since Hare's
computation~\cite{Hare2023} covers all candidate factors of degree $\le 15$ apart
from $975$ unpublished exceptions, the cases $k\le 15$ of Theorem~\ref{thm:main} may
be partly covered there; the theorem is certainly new for $k\ge 16$.
\end{remark}

\subsection*{Strategy}
We fix $a$ and $k$ and study all sequences $(q_n)$ compatible with the
recurrence $q_n=r_n-aq_{n-3}-q_{n-k}$, $r_n\in\{0,1\}$, $q_n\ge 0$
(Section~\ref{sec:setup}). They form a binary tree that branches only at indices
where $aq_{n-3}+q_{n-k}=0$. There is a distinguished \emph{main branch} whose
coefficients are given in closed form by an exact class formula, Lemma~\ref{lem:W}
(Section~\ref{sec:main}). Elementary lemmas show that every branch deviating from
the main branch dies (Section~\ref{sec:dev}), provided the main branch has a
positivity property at its forced indices and eventually dies. This is checked by
hand for $a\in(1/2,1)$ (Section~\ref{sec:large}) and by one interval-arithmetic
certificate valid for all $k\ge 67$ and $a\in(0,1/2]$ (Section~\ref{sec:U}). The
range $4\le k\le 149$ is covered by direct computer certificates
(Section~\ref{sec:small}).

\subsection*{Notation}
For a power series $G(a)=\sum_{i\ge 0} g_ia^i$ and an integer $L$ we write
$\Tr_L[G]=\sum_{i=0}^{L}g_ia^i$, with $\Tr_L[G]=0$ for $L<0$. Note
$(1+a)^{-(p+1)}=\sum_{i\ge 0}(-1)^i\binom{p+i}{p}a^i$. Throughout,
$c\in\{1,2\}$ denotes $k \bmod 3$ and residue classes are taken modulo $3$:
$\Cz=\{n\equiv 0\}$, $\Ck=\{n\equiv k\}$, $Z=\{n\equiv -k\}$; these are the three
distinct classes.

% ---------------------------------------------------------------------------
\section{The forced recurrence}\label{sec:setup}

\begin{definition}
Fix $a\in(0,1)$ and $k$. A \emph{path} is a sequence $r=(r_n)_{n\ge 0}\in\{0,1\}^{\mathbb{N}}$
with $r_0=1$. It determines $q_n=q_n(r)$ by $q_n=0$ for $n<0$ and
\begin{equation}\label{eq:rec}
q_n+aq_{n-3}+q_{n-k}=r_n\qquad(n\ge 0),
\end{equation}
equivalently $(1+ax^3+x^k)\sum q_nx^n=\sum r_nx^n$ as formal power series.
Put $S_n=aq_{n-3}+q_{n-k}$. The path \emph{dies at $N$} if $q_n\ge 0$ for
$n<N$ and $q_N<0$.
\end{definition}

\begin{lemma}\label{lem:reduce}
If $1+ax^3+x^k$ divides a $0$--$1$ polynomial with a nonzero cofactor with
nonnegative coefficients, then some path never dies.
\end{lemma}

\begin{proof}
Let $PQ=R$. Dividing $Q$ and $R$ by the largest power of $x$ dividing $Q$ we may
assume $q_0>0$. Then $r_0=q_0>0$, so $r_0=q_0=1$. The coefficient sequences of
$Q$ and $R$ (extended by zeros) satisfy~\eqref{eq:rec} with $q_n\ge 0$ for all $n$.
\end{proof}

Given $q_0,\dots,q_{n-1}\ge 0$, the choice $r_n=0$ keeps $q_n\ge 0$ iff $S_n=0$,
and $r_n=1$ keeps $q_n\ge 0$ iff $S_n\le 1$. In particular the paths that are
alive up to index $n$ branch only at indices with $S_n=0$; if $S_n>0$ then
$r_n=1$ is forced. By Lemma~\ref{lem:reduce}, Theorem~\ref{thm:main} follows once
we show that \emph{every path dies}.

\begin{lemma}[Linear response]\label{lem:resp}
Let $g_\Delta=[x^\Delta](1+ax^3+x^k)^{-1}$. For two paths $r,r'$,
$q_t(r')-q_t(r)=\sum_{s\le t}(r'_s-r_s)\,g_{t-s}$. For $0\le\Delta<k$ we have
$g_\Delta=(-a)^{\Delta/3}$ if $3\mid\Delta$ and $g_\Delta=0$ otherwise.
\end{lemma}

\begin{proof}
The first claim is linearity of $Q=R/P$. Writing
\begin{equation}\label{eq:pexp}
\frac{1}{1+ax^3+x^k}=\sum_{p\ge 0}(-1)^p\frac{x^{pk}}{(1+ax^3)^{p+1}},
\end{equation}
only $p=0$ contributes to degrees $<k$, and $(1+ax^3)^{-1}=\sum_i(-a)^ix^{3i}$.
\end{proof}

% ---------------------------------------------------------------------------
\section{The main branch and the class formula}\label{sec:main}

From now on $k\ge 4$ and $k\not\equiv 0\pmod 3$. Let $\F=Z\cap(k,2k]$ (the
\emph{free set}); its least element is $d_0=k+1$ if $c=1$ and $d_0=k+2$ if
$c=2$, and its largest element is $2k$.

\begin{definition}
The \emph{main branch} is the path $r^*$ with $r^*_n=[3\mid n]$ for $n<k$,
$r^*_k=1$, $r^*_n=0$ for $n\in\F$, $r^*_n=1$ for $n\in(k,2k]\setminus\F$, and
$r^*_n=1$ for $n>2k$. We write $q^*_n=q_n(r^*)$ and $S^*_n$ accordingly.
\end{definition}

On each residue class the support of $r^*$ is a tail: $\{n\in\Cz:n\ge 0\}$,
$\{n\in\Ck:n\ge k\}$ and $\{n\in Z:n\ge 2k+3\}$. Define the thresholds
$\tau=0,k,2k+3$ for $M$ in $\Cz,\Ck,Z$ respectively.

\begin{lemma}[Class formula]\label{lem:W}
For every $N\ge 0$, with $M_p=N-pk$ and $L_p=(M_p-\tau(M_p))/3$,
\[
q^*_N=\sum_{p\ge 0,\;M_p\ge 0}(-1)^p\,\Tr_{L_p}\!\big[(1+a)^{-(p+1)}\big],
\]
where terms with $L_p<0$ vanish. The right-hand side is a polynomial in $a$.
\end{lemma}

\begin{proof}
By~\eqref{eq:pexp}, $q^*_N=\sum_p(-1)^p\sum_{i\ge 0}\binom{p+i}{p}(-a)^i\,r^*_{M_p-3i}$.
All indices $M_p-3i$ lie in the class of $M_p$, and $r^*_{M_p-3i}=1$ iff
$M_p-3i\ge\tau(M_p)$, i.e.\ iff $i\le L_p$ (note $M_p\equiv\tau(M_p)$, so $L_p$ is an integer).
\end{proof}

The formula was checked against the recurrence with exact rational
polynomial arithmetic for $4\le k\le 120$ and all $N\le 5k+12$
($24999$ values, no discrepancy); see Section~\ref{sec:repro}.

\begin{corollary}\label{cor:C0}
For $t\in\Cz$ with $0\le t<3k$,
$q^*_t=\Tr_{t/3}[(1+a)^{-1}]=\dfrac{1-(-a)^{t/3+1}}{1+a}>0$ for all $a\in(0,1)$.
\end{corollary}

\begin{proof}
For $p=1$, $M_1=t-k\in Z$ and $M_1\ge 2k+3$ would need $t\ge 3k+3$; for $p=2$,
$M_2=t-2k\in\Ck$ and $M_2\ge k$ would need $t\ge 3k$; for $p\ge 3$, $M_p<0$.
\end{proof}

\begin{lemma}\label{lem:early}
For all $a\in(0,1)$: $q^*_n=0$ for $n<k$ with $n\notin\Cz$; $q^*_k=0$;
and for $n\in\F$ one has $q^*_{n-3}=q^*_{n-k}=0$, hence $S^*_n=0$ and $q^*_n=0$.
\end{lemma}

\begin{proof}
For $n<k$, \eqref{eq:rec} reads $q_n=r_n-aq_{n-3}$, so classes $\Ck$ and $Z$ stay
zero below $k$ under $r^*$. Then $q^*_k=1-aq^*_{k-3}-q^*_0=0$ since $k-3\in\Ck$.
For $n\in\F$, $n-3\in Z$ is either $<k$ or in $\F$ (induction), and
$n-k\in\Ck\cap(0,k]$, where $q^*$ vanishes.
\end{proof}

% ---------------------------------------------------------------------------
\section{Deviations from the main branch}\label{sec:dev}

\begin{lemma}[Window $n\le k$]\label{lem:window}
Let $a\in(0,1)$ and let $r$ be a path alive up to index $k$. Then
$r_{3j}=1$ and $q_{3j}=\sum_{i=0}^{j}(-a)^i\in(0,1]$ for $3j<k$;
$q_n=0$ for $n<k$, $n\in\Ck$; and $r_k=1$, $q_k=0$.
\end{lemma}

\begin{proof}
For $n<k$ we have $q_n=r_n-aq_{n-3}$. If $q_{n-3}\in(0,1]$ then $r_n=0$ would give
$q_n<0$, so $r_n=1$ and $q_n=1-aq_{n-3}\in[1-a,1)\subset(0,1]$. Starting from
$q_0=1$ this gives the claim on $\Cz$. If some $n<k$ in $\Ck$ had $q_n>0$, the same
propagation would give $q_{k-3}>0$ (as $k-3\in\Ck$), and then
$q_k=r_k-aq_{k-3}-1<0$. Hence $q_{k-3}=0$ and $q_k=r_k-1\ge 0$ forces $r_k=1$, $q_k=0$.
\end{proof}

So a path alive at index $k$ agrees with $r^*$ below $k$ except possibly on
$Z\cap[1,k)$, where $r^*$ vanishes.

\begin{lemma}[$\Cz$-shadowing]\label{lem:shadow}
Let $a\in(0,1)$, let $r$ agree with $r^*$ on $[0,d)$ for some $d\ge 1$, and let
$T\le\min(d+k-1,\,3k-1)$. If $r$ is alive up to $T$, then for every $t\in\Cz\cap[0,T]$
we have $r_t=1$ and $q_t(r)=q^*_t>0$.
\end{lemma}

\begin{proof}
Induction on $t\in\Cz$. For $t<d$ there is nothing to prove. Let $d\le t\le T$.
By induction (or Lemma~\ref{lem:window} if $t-3<d$) $q_{t-3}(r)=q^*_{t-3}>0$
(Corollary~\ref{cor:C0}, as $t-3<3k$), so $r_t=1$. By Lemma~\ref{lem:resp},
$q_t(r)-q^*_t=\sum_{s=d}^{t}(r_s-r^*_s)g_{t-s}$ with $0\le t-s\le T-d<k$; thus
$g_{t-s}\ne 0$ only for $s\in\Cz$, where $r_s=1=r^*_s$ by induction.
Positivity is Corollary~\ref{cor:C0}.
\end{proof}

\begin{proposition}[Late starts]\label{prop:late}
Let $a\in(0,1)$. If a path $r$ alive at $k$ has $r_{n_0}=1$ for some
$n_0\in Z\cap[1,k)$, then it dies at an index $\le k+n_0$.
\end{proposition}

\begin{proof}
Take $n_0$ minimal. Then $r$ agrees with $r^*$ on $[0,n_0)$ (Lemma~\ref{lem:window})
and $q_{n_0}(r)=1-aq_{n_0-3}(r)=1$, since $n_0-3\in Z$ lies below $n_0$.
Now $k+n_0\equiv k-k\equiv 0$, so $t=k+n_0-3\in\Cz$, and
$T=k+n_0-3\le\min(n_0+k-1,3k-1)$. If $r$ is alive up to $T$,
Lemma~\ref{lem:shadow} gives $q_t(r)=q^*_t>0$, and then
$q_{k+n_0}(r)=r_{k+n_0}-a\,q_t(r)-1<0$.
\end{proof}

\begin{proposition}[Free deviations]\label{prop:free}
Let $a\in(0,1)$ and let $r$ agree with $r^*$ on $[0,d_1)$ with $d_1\in\F$ and
$r_{d_1}=1$. Then $r$ dies at an index $\le d_1+k$.
\end{proposition}

\begin{proof}
By Lemma~\ref{lem:early}, $q_{d_1}(r)=1-aq^*_{d_1-3}-q^*_{d_1-k}=1$. The index
$t=d_1+k-3$ lies in $\Cz$ and $t\le 3k-3$, so with $T=t\le\min(d_1+k-1,3k-1)$
Lemma~\ref{lem:shadow} gives $q_t(r)=q^*_t>0$ if $r$ is alive up to $t$. Hence
$q_{d_1+k}(r)=r_{d_1+k}-aq_t(r)-1<0$.
\end{proof}

\begin{corollary}[Reduction to the main branch]\label{cor:reduce}
Let $a\in(0,1)$. Suppose the main branch dies at some $N>k$ and $S^*_n>0$ for
every $n\in(k,N)\setminus\F$. Then every path dies.
\end{corollary}

\begin{proof}
Let $r$ be a path alive at $k$ (otherwise we are done); by Lemma~\ref{lem:window},
$r$ agrees with $r^*$ on $[0,k]$ outside $Z\cap[1,k)$. If $r$ differs from $r^*$ on
$Z\cap[1,k)$, Proposition~\ref{prop:late} applies. Otherwise let $n>k$ be the first
index with $r_n\ne r^*_n$, if any. If $n\in\F$ then $r_n=1$ and
Proposition~\ref{prop:free} applies. If $n\notin\F$ and $n<N$, then $r_n=0$ while
$S_n(r)=S^*_n>0$, so $q_n(r)<0$. If no such $n<N$ exists, $r$ agrees with $r^*$ below
$N$; since $q^*_N=1-S^*_N<0$ forces $S^*_N>1$, both choices of $r_N$ give $q_N<0$.
\end{proof}

\begin{remark}\label{rem:trivial}
If $k\equiv 0\pmod 3$, the argument of Lemma~\ref{lem:window} gives $q_{k-3}>0$
in every path alive at $k-1$ (now $k-3\in\Cz$), so every path dies at $n=k$.
\end{remark}

% ---------------------------------------------------------------------------
\section{The range \texorpdfstring{$1/2<a<1$}{1/2 < a < 1}}\label{sec:large}

\begin{lemma}[Lemma A]\label{lem:A}
Let $k\ge 7$, $k\not\equiv 0\pmod 3$ and $a\in(1/2,1)$. Then
$q^*_{k+3}=a$ and $q^*_{k+6}=a(1-2a)<0$, and every path dies.
\end{lemma}

\begin{proof}
By Lemmas~\ref{lem:window} and~\ref{lem:early}, $q^*_k=0$, $q^*_3=1-a$, and
$q^*_6=1-a+a^2$ since $6<k$. Then $q^*_{k+3}=1-a\cdot 0-(1-a)=a$ and
$q^*_{k+6}=1-a^2-(1-a+a^2)=a(1-2a)<0$. The indices $n\in(k,k+6)\setminus\F$ are
$k+3\in\Ck$, with $S^*_{k+3}\ge q^*_3>0$, and the $\Cz$-indices among
$k+1,\dots,k+5$, with $S^*_n\ge aq^*_{n-3}>0$ by Corollary~\ref{cor:C0}. Hence the
main branch dies at some $N\le k+6$ and $S^*_n>0$ on $(k,N)\setminus\F$;
Corollary~\ref{cor:reduce} concludes.
\end{proof}

\begin{remark}
One can also see directly that every path whose first deviation from $r^*$ is at
an index $s>k$ dies at $k+6$: the free indices below $k+6$ are $d_0$ and $d_0+3$,
so the relevant differences $k+3-s$ and $k+6-s$ lie in $\{1,2,4,5\}$ and the
corresponding $g_\Delta$ vanish by Lemma~\ref{lem:resp}. Thus $q_{k+6}\le a(1-2a)<0$
whatever $r_{k+6}$ is.
\end{remark}

\begin{remark}[Formal verification]\label{rem:lean}
Three of the hand arguments have been checked in Lean~4 (v4.34.1, with Mathlib),
for an arbitrary path $r\in\{0,1\}^{\mathbb{N}}$ and without coprimality
hypotheses; the theorems are stated for general $b$ and are used here with $b=3$
(file \texttt{lean/Trinom/Trinom/Ogish.lean}): \texttt{A\_b\_star} gives the
conclusion ``every path dies'' of Lemma~\ref{lem:A} (no path has $q_n\ge 0$ for all
$n\le k+6$); \texttt{E1} gives the step of Lemma~\ref{lem:window} that a positive
coefficient in $\Ck$ below $k$ forces $q_k<0$; and \texttt{dies\_of\_dvd} gives
Remark~\ref{rem:trivial}. The file contains no \texttt{sorry}, and
\texttt{\#print axioms} reports only \texttt{propext}, \texttt{Classical.choice} and
\texttt{Quot.sound}. The remaining lemmas and the computer certificates are not
formalised.
\end{remark}

% ---------------------------------------------------------------------------
\section{The range \texorpdfstring{$0<a\le 1/2$, $k\ge 67$}{0 < a <= 1/2, k >= 67}: a uniform certificate}\label{sec:U}

\begin{proposition}[Lemma U]\label{prop:U}
Let $k\ge 67$, $k\not\equiv 0\pmod 3$ and $a\in(0,1/2]$. Then the main branch
dies at some $N\le 5k+12$, and $S^*_n>0$ for every $n\in[k,N)\setminus\F$.
\end{proposition}

The proof is a rigorous computation that is independent of $k$; we first explain
why finitely many computations suffice.

\subsection{Reduction to finitely many types}
Write $N=jk+m$ with $0\le m<k$. Then $M_p=(j-p)k+m$ lies in the class
$(j-p)c+m\bmod 3$, and Lemma~\ref{lem:W} gives
$L_p=(\alpha_pk+m-\beta_p)/3$, where $(\alpha_p,\beta_p)=(j-p,0)$, $(j-p-1,0)$ or
$(j-p-2,3)$ for $M_p$ in $\Cz$, $\Ck$, $Z$ respectively. Terms with $\alpha_p<0$
vanish. Terms with $\alpha_p=0$ (``live'' terms) have $L_p=(m-\beta_p)/3$,
independent of $k$. Terms with $\alpha_p\ge 1$ (``small'' terms) have
$L_p\ge(k-3)/3$.

\begin{lemma}\label{lem:tail}
Let $p\ge 0$, $L\ge p$ and $a\in(0,1/2]$. Then
$\big|(1+a)^{-(p+1)}-\Tr_L[(1+a)^{-(p+1)}]\big|\le\binom{p+L+1}{p}a^{L+1}$, and the
right-hand side is non-increasing in $L$.
\end{lemma}

\begin{proof}
The series $\sum_i(-1)^i\binom{p+i}{p}a^i$ alternates, and the ratio of consecutive
terms is $\frac{p+i+1}{i+1}a\le\frac{p+i+1}{2(i+1)}\le 1$ for $i\ge p-1$; apply the
Leibniz bound. The same ratio computation with $i=L+1$ gives monotonicity.
\end{proof}

Consequently, for $j\le 5$ and $k\ge K$,
\[
\begin{gathered}
q^*_{jk+m}(a)=\frac{\mathcal{N}^{(c)}_{j,m}(a)}{(1+a)^6}+E,\\
|E|\le\sum_{p:\,\alpha_p\ge 1}\binom{p+\Lambda+1}{p}a^{\Lambda+1},\qquad
\Lambda=\lfloor (K-3)/3\rfloor,
\end{gathered}
\]
where $\mathcal{N}^{(c)}_{j,m}$ is an explicit polynomial: each small term
contributes $(-1)^p(1+a)^{5-p}$ and each live term $(-1)^p\Tr_{L_p}[(1+a)^{-(p+1)}](1+a)^6$.
It depends on $k$ only through $c=k\bmod 3$. For $m>M_0$ (we use $M_0=60$) the live
terms are also replaced by $(1+a)^{-(p+1)}$, with error at most
$\binom{p+\Lambda_0+1}{p}a^{\Lambda_0+1}$, $\Lambda_0=19$; the resulting type depends
on $m$ only through $m\bmod 3$. Hence, for fixed $c$, the main branch on the
windows $j=1,\dots,5$ is described by finitely many \emph{types}, each a rational
function plus an explicit error bound.

\subsection{The certificate}
The script \texttt{retsenziya/uniform\_sertifikat.py} does the following for each
$c\in\{1,2\}$ and each of the $64$ intervals $I=[i/128,(i+1)/128]$, $0\le i\le 63$.
It walks the main branch through the positions $(j,m)$ in increasing order:
$j=1,\dots,4$ with $m=0,\dots,M_0+3$ individually followed by the three residues of
$m>M_0+3$ (``generic'' positions), and $j=5$ with $m=0,\dots,12$; this covers every
$n\in[k,5k+12]$ when $k>M_0+3$. For each type it decides the sign on $I$ rigorously:
after removing the largest power $a^v$ dividing $\mathcal{N}$, it encloses
$\mathcal{N}(a)/a^v$ on $I$ by Arb ball arithmetic (\texttt{arb\_poly} evaluated on
the ball $[lo,hi]$) and compares with the enclosure
$(1+hi)^6\sum C\,hi^{\lambda-v}$ of the error term; a sign is accepted only if it is
strict on the whole closed interval, otherwise the type is treated as ambiguous
(possibly zero). Then:
\begin{itemize}
\item[(S1)] at every non-free position before the death, $q^*_{n-3}$ or $q^*_{n-k}$
must be certified strictly positive on $I$;
\item[(S2)] some position must be certified strictly negative on $I$ (the death).
\end{itemize}
Failing intervals would be bisected (up to depth $40$); none had to be.
Since the walk is in increasing order, (S1) is checked at every non-free $n$ up to the
first certified negative value, which bounds the actual death index from above.
Because intervals are closed and signs are strict on them, isolated parameters
(e.g.\ roots of some $q^*_n$) need no separate treatment.

The certificate was run with $K=151$ and with $K=67$ (the condition $k>M_0+3=63$
holds). In both runs and in both classes all $64$ intervals were certified without
subdivision, the deaths being of the four types
\[
3k+12\ (2\text{ intervals}),\quad 4k+15\ (1),\quad 5k+3\ (59),\quad 5k+12\ (2).
\]
A run takes about $0.7$ seconds. This proves Proposition~\ref{prop:U}, the interval
$[63/128,1/2]$ covering $a=1/2$.

\begin{remark}[Sanity checks]
Three deliberate mutations of the certificate fail as they should: checking (S1)
also on the free set makes all $128$ intervals fail; flipping the sign of the live
terms changes the outcome; inflating the error bound by $10^{12}$ makes $44$
intervals fail. A comparison of the type signs with exactly computed $q^*_N$ at
$1200$ random parameters for $k=151,152,301$ found no contradiction; this comparison
uses the same type model and is therefore not an independent check.
\end{remark}

\begin{remark}[Limit shapes; not used in the proof]\label{rem:limit}
Lemma~\ref{lem:W} explains the four death types. For instance, for $k\ge 13$,
\[
q^*_{5k+3}=\Tr_{k+1}[G_1]+(-1)^ka^{k+1},\qquad
q^*_{5k+12}=\Tr_{k+4}[G_2]+(-1)^{k+1}a^{k+4},
\]
with $G_1=a-1+(1+a)^{-1}-(1+a)^{-2}+(1+a)^{-3}=a(a^3+2a^2+a-1)/(1+a)^3$ and
$G_2=-56a^4+41a^3-16a^2+5a-1+(1+a)^{-1}-(1+a)^{-2}+(1+a)^{-3}$
(verified by computer for $13\le k\le 300$). Numerically (coefficients compared up to
$k=40$), the $k$-independent limits of $q^*_n$ at $n=5k+3,5k+12,4k+15,3k+12$ are
negative for $a<\alpha$, $a>\beta$, $a>0.47098\ldots$ and $a>0.47681\ldots$
respectively, where $\alpha=0.4655712318\ldots$ and $\beta=0.4484743767\ldots$
are the roots in $(0,1)$ of $1-a-2a^2-a^3$ and of the numerator of $G_2$. Since
$(0,\alpha)\cup(\beta,1/2]=(0,1/2]$, the main branch always dies; the certificate
makes this rigorous uniformly in $k$.
\end{remark}

% ---------------------------------------------------------------------------
\section{Small \texorpdfstring{$k$}{k} and isolated parameters}\label{sec:small}

\begin{proposition}\label{prop:small}
For $4\le k\le 149$, $k\not\equiv 0\pmod 3$, and every $a\in(0,1)$, every path dies.
\end{proposition}

This is established by an exhaustive computer certificate over the whole parameter
interval $(0,1)$ (script \texttt{simvolik\_isbot\_flint.py} with
\texttt{GCD\_TEKSHIR=1}). The program performs a depth-first search over paths in
which each node carries a subinterval $I\subset(0,1)$ and the exact polynomials
$q_0,\dots,q_{n-1}\in\mathbb{Q}[a]$ valid on $I$ (FLINT \texttt{fmpq\_poly}).
At index $n$:
\begin{itemize}
\item the branch $r_n=0$ is opened when $q_{n-3}$ and $q_{n-k}$ are both the zero
polynomial;
\item for $r_n=1$, $I$ is split according to the sign of $q_n=1-aq_{n-3}-q_{n-k}$.
Real roots are counted exactly with Sturm sequences (PARI \texttt{polsturm});
approximate roots from \texttt{polrootsreal} are enclosed in rational boxes of
width about $2^{-40}$ whose root counts are confirmed by \texttt{polsturm}, with
exact bisection as a fallback. An Arb ball evaluation that excludes $0$ on $I$
certifies a constant sign without Sturm. Subintervals where $q_n<0$ die; the others,
including the narrow boxes containing roots, are continued (conservatively);
\item parameters where $r_n=0$ is possible without $q_{n-3},q_{n-k}$ vanishing
identically are isolated algebraic numbers: common roots of $q_{n-3}$ and $q_{n-k}$
(gcd, factored by PARI \texttt{factor}) and roots of $q_n$ falling exactly on a
subinterval endpoint are collected.
\end{itemize}
At each collected algebraic parameter $\alpha$ the complete tree, including all
$r_n=0$ branches, is rerun in exact arithmetic in $\mathbb{Q}(\alpha)$
(\texttt{nuqtaviy.py}); signs are decided by refining an isolating interval of the
minimal polynomial with Sturm sequences. A run is accepted if no branch reaches the
depth bound $N$, no terminating cofactor is detected, and all isolated parameters die.

For $k\le 60$ the runs (depth bound $N=2000$) are recorded in
\texttt{retsenziya/toliq\_census.log}; all $38$ runs are accepted, $36$ isolated
parameters were collected in total and all die, and for $k\ge 13$ the maximal depth
reached is exactly $5k+12$. For $61\le k\le 149$ ($60$ values,
$N=4000$; \texttt{retsenziya/b3\_61\_150.log}) all $60$ runs are accepted, the maximal
depth is exactly $5k+12$ in each run, $60$ isolated parameters were collected in total
(for instance $a=1/2$) and all of them die; each run takes at most about $7$ seconds.
The range $67\le k\le 149$ is thus covered twice (here and by
Proposition~\ref{prop:U}).

\subsection*{Exact check at root points for \texorpdfstring{$151\le k\le 250$}{151 <= k <= 250}}
As an independent confirmation of the treatment of isolated parameters in
Proposition~\ref{prop:U}, for each $k\in[151,250]$, $k\not\equiv 0$ ($67$ values), the
script \texttt{retsenziya/ildiz\_nuqta.py} computes all main-branch polynomials
$q^*_n$, $n\le 5k+12$, collects all their roots in $(0,1/2]$, and runs the complete
exact tree in $\mathbb{Q}(\alpha)$ at each root (depth bound $5k+60$). In every case
there are exactly six such roots, the same for all $k$:
\[
\begin{gathered}
0.4484743767\ldots,\quad 0.4655712319\ldots,\quad 0.4709761380\ldots,\\
0.4768143659\ldots,\quad 0.4874322956\ldots,\quad 1/2
\end{gathered}
\]
(the last being the root of $q^*_{k+6}=a(1-2a)$). All $402$ point trees die, none
survives, and the maximal depth is exactly $5k+12$ for every $k$; in total
$1895851$ branchings were explored (about $19971$ s of running time summed over the
runs). With the depth bound lowered to $5k+10$, two points are correctly reported
alive. An independent re-implementation produced the same outcome for
$k=151,152,154,155,301,302$.

% ---------------------------------------------------------------------------
\section{Proof of Theorem~\ref{thm:main}}

Let $k\ge 4$, $k\not\equiv 0\pmod 3$, $a\in(0,1)$. By Lemma~\ref{lem:reduce} it
suffices to show that every path dies. If $k\le 149$ this is
Proposition~\ref{prop:small}. If $k\ge 67$ and $a\in(1/2,1)$, it
is Lemma~\ref{lem:A}. If $k\ge 67$ and $a\in(0,1/2]$, Proposition~\ref{prop:U} gives
the hypotheses of Corollary~\ref{cor:reduce} (note $n=k\notin\F$ and
$S^*_k=aq^*_{k-3}+q^*_0=1>0$), which concludes. \qed

The proof depends on the correctness of python-flint~\cite{pythonflint}
(FLINT~\cite{FLINT} and its ball-arithmetic component Arb~\cite{Arb}:
\texttt{fmpq\_poly}, \texttt{arb}, \texttt{arb\_poly}) and of
PARI/GP~\cite{PARI} (\texttt{polsturm}, \texttt{polrootsreal}, \texttt{factor}),
both standard libraries.

% ---------------------------------------------------------------------------
\section{Computer verification and reproducibility}\label{sec:repro}

\begingroup\raggedright
All scripts are in the project directory (paths relative to it) and were run with a
Python environment providing python-flint~0.9.0 and cypari2~2.2.4 (PARI~2.17.2),
denoted \texttt{python} below.
\begin{enumerate}
\item \emph{Class formula} (Lemma~\ref{lem:W}):
\texttt{python retsenziya/oyna\_formula.py 4 120} compares the formula with the
recurrence for all $N\le 5k+12$; output: $24999$ values, $0$ discrepancies.
\item \emph{Uniform certificate} (Proposition~\ref{prop:U}):
\texttt{python retsenziya/uniform\_sertifikat.py 151 60} and
\texttt{python retsenziya/uniform\_sertifikat.py 67 60}; output per class:
$64$ intervals checked, $0$ failed, death types
$\{(3,12)\!:\!2,\ (4,15)\!:\!1,\ (5,3)\!:\!59,\ (5,12)\!:\!2\}$, where $(j,m)$ stands for
$n=jk+m$.
\item \emph{Small $k$} (Proposition~\ref{prop:small}):
\texttt{GCD\_TEKSHIR=1 python simvolik\_isbot\_flint.py 3 $k$ 2000} for $k\le 60$
(driver \texttt{retsenziya/toliq\_census.sh}, log \texttt{retsenziya/toliq\_census.log}) and
\texttt{GCD\_TEKSHIR=1 python simvolik\_isbot\_flint.py 3 $k$ 4000} for $61\le k\le 149$
(driver \texttt{retsenziya/b3\_61\_150.sh}, log \texttt{retsenziya/b3\_61\_150.log}). Each run
prints a JSON record; acceptance is the field \texttt{isbot\_toliq}.
\item \emph{Root points}: \texttt{python retsenziya/ildiz\_nuqta.py $k$} for
$151\le k\le 250$ (driver \texttt{retsenziya/ildiz\_nuqta.sh}, log
\texttt{retsenziya/ildiz\_nuqta.log}).
\item \emph{Formal verification} (Remark~\ref{rem:lean}): \texttt{lake build} in
\texttt{lean/Trinom} (toolchain \texttt{leanprover/lean4:v4.34.1}, Mathlib
\texttt{v4.34.1}); the build completes without errors and the files
\texttt{Trinom/Basic.lean} and \texttt{Trinom/Ogish.lean} contain no \texttt{sorry}.
\end{enumerate}

\subsection*{Availability}
The scripts listed above (\texttt{simvolik\_isbot\_flint.py}, \texttt{nuqtaviy.py} and
the files under \texttt{retsenziya/}), the complete logs of all runs reported in
Sections~\ref{sec:U} and~\ref{sec:small}, and the Lean sources are
available together with this paper. Each certificate run recomputes its verdict and
reads no stored data (the scripts only write their output), so the logs serve only
for comparison. Timings quoted in the text are wall-clock times of single
runs and are given for orientation only.
All of this material is deposited at Zenodo, doi:\,\ZenodoDOI,
and mirrored at \url{https://github.com/selloriybrendi/zero-one-trinomial-b3}.
All runs were made on an Apple M5 (10 cores, 16\,GB RAM) under macOS, with
python-flint~0.9.0 and PARI/GP~2.17.2.
\par\endgroup

% ---------------------------------------------------------------------------
\section*{Use of AI tools}
The AI assistant Claude (Anthropic) was used in the preparation of this work: to
write and run the computational scripts, including the certificate programs of
Sections~\ref{sec:U} and~\ref{sec:small}, to draft the Lean formalisation of
Remark~\ref{rem:lean}, to search and check the literature, and to draft the text of
this paper. The mathematical claims do not rest on the output of the AI tool: the
computer-assisted parts are checked by the exact and interval arithmetic of
FLINT/Arb~\cite{FLINT,Arb} and PARI/GP~\cite{PARI}, and the formalised parts by the
Lean proof checker. The author has checked all arguments, computations and
references in this paper and takes full responsibility for its content.

% ---------------------------------------------------------------------------
\section{Outlook}

For general odd $b$ the class formula extends: if $\gcd(b,k)=1$ and the main branch
is defined by $r_n=1$ unless $q_{n-b}$ and $q_{n-k}$ both vanish identically, the
thresholds are $\tau_0=0$, $\tau_1=k$, $\tau_j=jk+b$ ($j\ge 2$) on the class
$n\equiv jk\pmod b$, and $q_N=\sum_{p:\,L_p\ge 0}(-1)^p\Tr_{L_p}[(1+a)^{-(p+1)}]$ with
$L_p=(M_p-\tau_{j(M_p)})/b$. This was checked by computer (no discrepancy) for
$b=3$ ($4\le k\le 60$), $b=5$ ($6\le k\le 40$) and $b=7$ ($8\le k\le 36$), but we do
not state it here as a proved lemma for $b\ge 5$. The following is known only
partially and is \emph{not} proved here.

\begin{itemize}
\item \emph{Finite ranges.} Exhaustive certificates of the type of
Section~\ref{sec:small} cover $b=5$, $6\le k\le 39$, and $b=7$, $8\le k\le 30$
($k\not\equiv 0\pmod b$).
\item \emph{Main branch, uniformly in $k$.} The analogue of the certificate of
Section~\ref{sec:U} shows that the main branch dies by index $7k+5$ for $b=5$, all
$k\ge 201$, and by $9k+7$ for $b=7$, all $k\ge 1001$; the dominant death types are
$(b+2)k+b$ for $b=3,5,7$. For $a\in(1/2,1)$ and $k>2b$ the main branch has
$q_{k+b}=a$ and $q_{k+2b}=a(1-2a)<0$.
\item \emph{Deviations.} The free-deviation lemma (Proposition~\ref{prop:free}) fails
for $b\ge 5$: for $b=5$, $k=31$, a deviation at $n=2k$ survives up to $13k+30$ at an
isolated $a\approx 0.2469$. A uniform treatment of deviating branches is missing.
\item \emph{Stability of deviation patterns.} For $b=5$, describing the leaves of
the full tree in window coordinates (offsets $\le 30$ from window ends kept exactly,
middle offsets only modulo $b$), the sets of leaf patterns for $k=66$ and $k=71$
contain $53151$ and $54617$ patterns, with $53147$ in common; only $4$ patterns occur
for $k=66$ alone. The trees have $66706$ and $85329$ leaves, no live branch, and at
most four deviations per path.
\end{itemize}

\begin{conjecture}
For $b=5$ (and plausibly every odd $b$), the deviation tree stabilises in window
coordinates for large $k$, so that a finite catalogue of patterns, each with a uniform
interval certificate, would prove the analogue of Theorem~\ref{thm:main}.
\end{conjecture}

\raggedright
\bibliographystyle{amsplain}
\bibliography{b3_trinom}

\end{document}
